% Diese Seite verwendet eine eigene Seitengeometrie und verzichtet auf die Fußnote.
\newgeometry{left=2cm, right=2cm, top=1cm, bottom=1cm}
\thispagestyle{empty}

% Überschrift der Seite, zentriert
\centering
\section*{\myul{Terminologie}}

\normalsize

\vspace{-4mm}

\begin{table}[H]
    \raggedright
    \begin{tabular}{l l l}
        \textbf{Begriff} & \textbf{Bedeutung} & \textbf{Mathematisch} \\[2.0ex]
        Stromlinie   & \begin{tabular}[c]{@{}l@{}}Linie, die zu jedem Zeitpunkt tangential zum örtlichen\\ Geschwindigkeitsvektor verläuft.\end{tabular} & $\varPsi =$ konstant \\[3.0ex]
        Stromfaden   & \begin{tabular}[c]{@{}l@{}}Die Gesamtheit aller Stromlinien, die durch eine Fläche $S$ hindurch\\ treten (besteht aus der Ein- und Austrittsfläche $S_1$, $S_2$ und der\\ Mantelfläche).\end{tabular} & - \\[4.5ex]
        Stromröhre   & Bezeichnung für die röhrenförmige Mantelfläche des Stromfadens. & - \\[2.0ex]
        Bahnlinie    & Die Wegkurve eines einzelnen, individuellen Masseteilchens. & - \\[2.0ex]
        Streichlinie & \begin{tabular}[c]{@{}l@{}}Der Ort aller Massepunkte in einer Strömung, die an einem be-\\stimmten, festen Punkt vorbei gestrichen sind.\end{tabular} & - \\[3.0ex]
        Zeitlinie    & \begin{tabular}[c]{@{}l@{}}Linien, welche durch Teilchen- / Blasenbildung gleichzeitig entlang\\ eines gewählten Strömungspotenzials erzeugt werden.\end{tabular} & - \\[3.0ex]
        Isotache     & Linie gleicher Geschwindigkeit in der Strömung. & $\left|\underline{v}\right| =$ konstant \\[2.0ex]
        Isokline     & Linie gleicher Stromlinienneigung in der Strömung. & $\tan(\alpha) = \dfrac{v}{u} =$ konst. \\
    \end{tabular}
\end{table}

\vspace{2mm}

\centering
\section*{\myul{Symbolverzeichnis}}

\vspace{-4mm}

\begin{table}[H]
    \raggedright
    \begin{tabular}{c c l}
        \xrowht[()]{3.6mm} \textbf{Symbol}                              & \textbf{SI-Einheit}                          & \textbf{Bedeutung}                                  \\[2.0ex]
        \xrowht[()]{3.6mm} $c_\text{p}$                                 & $\left[-\right]$                             & Druckbeiwert                                        \\
        \xrowht[()]{3.6mm} $c_p$, $c_v$                                 & $\left[\frac{\text{J}}{\text{kg\,K}}\right]$ & Spezifische, isobare / isochore Wärmekapazität      \\
        \xrowht[()]{3.6mm} $c_\text{A}$, $c_\text{W}$                   & $\left[-\right]$                             & Auftriebs- und Widerstandsbeiwert                   \\
        \xrowht[()]{3.6mm} $F_x$, $F_y$, $F_z$                          & $\left[\text{N}\right]$                      & Kraft auf ein Fluid (in $x$, $y$ oder $z$-Richtung) \\
        \xrowht[()]{3.6mm} $u$, $v$, $w$                                & $\left[\frac{\text{m}}{\text{s}}\right]$     & Geschwindigkeiten in $x$, $y$ und $z$-Richtung      \\
        \xrowht[()]{3.6mm} $U_\infty$ / $v_\infty$                      & $\left[\frac{\text{m}}{\text{s}}\right]$     & Geschwindigkeit in $x$-Richtung, unendlich entfernt \\[3.0ex]
        \xrowht[()]{3.6mm} $\dot{\gamma}$                               & $\left[\frac{1}{\text{s}}\right]$            & Schergeschwindigkeit                                \\
        \xrowht[()]{3.6mm} $\zeta$                                      & $\left[-\right]$                             & Druckverlustbeiwert                                 \\
        \xrowht[()]{3.6mm} $\eta$                                       & $\left[\text{Pa\,s}\right]$                  & Kinematische Viskosität                             \\
        \xrowht[()]{3.6mm} $\vartheta$                                  & $\left[-\right]$                             & Umlenkwinkel (Bezogen auf Anströmrichtung)          \\
        \xrowht[()]{3.6mm} $\kappa$                                     & $\left[-\right]$                             & Isentropenexponent                                  \\
        \xrowht[()]{3.6mm} $\nu$                                        & $\left[\frac{\text{m}^2}{\text{s}}\right]$   & Dynamische Viskosität (Zähigkeit)                   \\
        \xrowht[()]{3.6mm} $\tau$ / $\tau_\text{W}$                     & $\left[\text{Pa}\right]$                     & Schubspannung / Wandschubspannung                   \\[3.0ex]
        \xrowht[()]{3.6mm} $\underline{x}$, $\underline{\underline{x}}$ &                                              & Allgemein: Vektor / Matrix                          \\
        \xrowht[()]{3.6mm} $\overline{x}$                               &                                              & Allgemein: Gemittelte Größe                         \\[3.0ex]
        \xrowht[()]{3.6mm} $\nabla(\,\cdots)$                           &                                              & Nabla-Operator (Mathematisch)                       \\
        \xrowht[()]{3.6mm} $\Delta(\,\cdots)$                           &                                              & Gradient (Mathematisch)                             \\
        \xrowht[()]{3.6mm} rot$(\,\cdots)$                              &                                              & Rotation (Mathematisch)                             \\
        \xrowht[()]{3.6mm} div$(\,\cdots)$                              &                                              & Divergenz (Mathematisch)                            \\
    \end{tabular}
\end{table}

{
    \vspace{1mm}
    \footnotesize
    Der Übersicht wegen sind hier nicht alle Symbole aufgeführt.\\
    Es kann zu doppelter Symbolbelegung kommen.
}


\restoregeometry